\documentclass[11pt,reqno]{amsart}
\usepackage{amsmath,amssymb,amsthm,mathtools}
\usepackage[margin=1.1in]{geometry}
\usepackage[T1]{fontenc}
\usepackage{lmodern}
\usepackage{hyperref}
\hypersetup{
  colorlinks=true,citecolor=blue,linkcolor=blue,urlcolor=blue,
  pdftitle={A negative answer to a question of Bourgain on L1 norms of exponential sums},
  pdfauthor={Paata Ivanisvili},
  pdfsubject={Exponential sums, L1 norms, flat polynomials},
  pdfkeywords={exponential sums, L1 norm, flat polynomials, Littlewood polynomials, Bourgain}
}

\newtheorem{theorem}{Theorem}[section]
\newtheorem{lemma}[theorem]{Lemma}
\newtheorem{proposition}[theorem]{Proposition}
\newtheorem{corollary}[theorem]{Corollary}
\theoremstyle{definition}
\newtheorem{definition}[theorem]{Definition}
\newtheorem{remark}[theorem]{Remark}

\newcommand{\T}{\mathbb{T}}
\newcommand{\Z}{\mathbb{Z}}
\newcommand{\R}{\mathbb{R}}
\newcommand{\C}{\mathbb{C}}
\newcommand{\e}{\mathrm{e}}
\newcommand{\cS}{\mathcal{S}}
\newcommand{\cD}{\mathcal{D}}

\begin{document}

\title[A negative answer to a question of Bourgain]{A negative answer to a question of Bourgain on $L^1$ norms of exponential sums}

\author{Paata Ivanisvili}
\address{Department of Mathematics, University of California, Irvine, CA 92697, USA}
\email{pivanisv@uci.edu}

\subjclass[2020]{42A05, 42A16, 11L03, 30C10}
\keywords{Exponential sums, $L^1$ norm, flat polynomials, Littlewood polynomials}

\begin{abstract}
Bourgain asked whether there is an absolute constant $\varepsilon>0$ such that
\[
\Big\|\,|S|^{-1/2}\sum_{s\in S}e^{2\pi is\theta}\Big\|_{L^1(\mathbb T)}\le1-\varepsilon
\]
for every finite set $S\subset\mathbb Z$ with $|S|\ge2$. We show that the answer is negative: the supremum of the left-hand side over all sets with $|S|=n$ tends to $1$ as $n\to\infty$. For every $\rho\in(0,1)$ we construct sets $S\subset\{0,1,\dots,N-1\}$ such that the balanced indicator $1_S-\rho$ satisfies $\max_\theta\big|\sum_{k<N}(1_S(k)-\rho)e^{2\pi ik\theta}\big|\le(1+o(1))\sqrt{\rho(1-\rho)N}$. The construction adapts to a fixed bias the recent construction of real Littlewood polynomials with asymptotically minimal maximum modulus. For these sets the normalized $L^1$ norm tends to $\sqrt{1-\rho}$, which is the largest possible value for subsets of an interval of density $\rho$.
\end{abstract}

\maketitle

\section{Introduction}

\subsection{Bourgain's question}
For $x\in\R$ write $\e(x)=e^{2\pi ix}$, and let $\T=\R/\Z$ carry the Lebesgue probability measure. For a finite nonempty set $S\subset\Z$ with $n=|S|$ put
\[
P_S(\theta)=\sum_{s\in S}\e(s\theta),\qquad X_S=n^{-1/2}P_S .
\]
Then $\|X_S\|_{L^2(\T)}=1$, so $\|X_S\|_{L^1(\T)}\le1$, with strict inequality if $n\ge2$, since then $|X_S|$ is not constant. Let
\[
\beta_n=\sup_{|S|=n}\|X_S\|_{L^1(\T)},\qquad \beta=\sup_{n\ge2}\beta_n .
\]
In his study of the spectral type of Ornstein's class one transformations \cite{Orn72}, whose spectral measures are generalized Riesz products of normalized $0$--$1$ polynomials, Bourgain \cite{Bou93} asked whether there is an absolute $\varepsilon>0$ such that
\begin{equation}\label{eq:bourgain}
\|X_S\|_{L^1(\T)}\le1-\varepsilon\qquad\text{for every finite }S\subset\Z,\ |S|\ge2,
\end{equation}
that is, whether $\beta<1$. Using the theorem of McGehee, Pigno and Smith, he proved \cite[Proposition~2]{Bou93} that $\beta_n\le1-c\log n/n$ for an absolute $c>0$. Thus $\beta_n<1$ for each $n$, and \eqref{eq:bourgain} is equivalent to $\limsup_n\beta_n<1$. As Aistleitner notes \cite{Ais13}, deciding whether $\beta<1$ would be a major step towards confirming or disproving the conjecture of Klemes and Reinhold \cite{KR97} that every rank one transformation has singular spectral type; see also \cite{Kle96,AH12}. Theorem~\ref{thm:negative} below shows that the uniform bound \eqref{eq:bourgain} is not available for that purpose; it does not by itself decide the conjecture.

Lower bounds for $\limsup_n\beta_n$ have been obtained by several authors. Karatsuba \cite{Kar98} proved $\|X_S\|_1\ge n/\sqrt{K(S)}$, where $K(S)$ is the number of solutions of $s_1+s_2=s_3+s_4$ in $S$; for Sidon sets this gives $\limsup_n\beta_n\ge1/\sqrt2$. Random subsets of density $\frac12$ give only $\sqrt{\pi/8}\approx0.63$, as observed in \cite{CE15} using \cite{BL01}. Aistleitner \cite{Ais13} proved $\limsup_n\beta_n\ge\sqrt\pi/2\approx0.886$, using a quantitative central limit theorem for lacunary sequences, and Choi and Erd\'elyi \cite{CE15} extended this to $L^p$ norms, $0<p<2$. Ivanisvili and Stone \cite{IS24} proved a quantitative form of the inequality $\|X\|_1\le\|X\|_2$ for random variables with $\|X\|_2=1$, in terms of $\|X\|_p$ and $\|X\|_q$ with $2<p<q$; they used it to bound $1-\|X_S\|_1$ from below for the squares $S=\{1^2,\dots,m^2\}$, and noted that \eqref{eq:bourgain} was open even for that family. Recently Sanders \cite{San26} proved $\limsup_n\beta_n\ge0.953\ldots$ by lifting Ruzsa's dense Sidon sets from a cyclic group to the integers; he also showed that lifts of near-perfect difference sets have normalized $L^1$ norms bounded away from $1$. Several of these results are surveyed by Erd\'elyi \cite{Erd20}.

\subsection{\texorpdfstring{$L^1$}{L1} norms of exponential sums}
Bourgain's question concerns the upper end of the range of $\|X_S\|_1$. The lower end is the subject of Littlewood's conjecture, posed by Hardy and Littlewood \cite{HL48}: there is an absolute $c>0$ such that $\|P_S\|_{L^1(\T)}\ge c\log|S|$ for every finite $S\subset\Z$ with $|S|\ge2$. It was proved independently by Konyagin \cite{Kon81} and by McGehee, Pigno and Smith \cite{MPS81}. In terms of the autocorrelations $r(h)=\#\{(a,b)\in S^2:a-b=h\}$ one has $|X_S|^2-1=n^{-1}\sum_{h\ne0}r(h)\e(h\theta)$, and \eqref{eq:autocorr} below shows that \eqref{eq:bourgain} is equivalent to the linear lower bound $\|\sum_{h\ne0}r(h)\e(h\,\cdot)\|_1\ge c\,|S|$, whereas the McGehee--Pigno--Smith theorem gives $c\log|S|$. Our result shows that the linear bound fails.

\subsection{Flat polynomials}
The negative answer requires polynomials with coefficients in $\{0,1\}$ whose modulus is nearly constant in $L^1$, and so it belongs to the circle of problems on flat polynomials. Littlewood \cite{Lit66,Lit68} asked for polynomials $P(z)=\sum_{k<N}\varepsilon_kz^k$ with $\varepsilon_k\in\{-1,1\}$ and $c_1\sqrt N\le|P(z)|\le c_2\sqrt N$ on $|z|=1$. The Rudin--Shapiro polynomials \cite{Sha51,Rud59} satisfy the upper bound with $c_2=\sqrt2$ when $N$ is a power of $2$. Erd\H{o}s \cite{Erd57} asked whether there is an absolute $c>0$ such that $\max_{|z|=1}|P(z)|\ge(1+c)\sqrt N$ for every polynomial of degree $N-1$ whose coefficients have modulus one. For complex coefficients Kahane \cite{Kah80} answered this negatively, constructing ultraflat polynomials with $(1-o(1))\sqrt N\le|P|\le(1+o(1))\sqrt N$ on the unit circle; earlier K\"orner \cite{Kor80} had constructed flat ones, and Bombieri and Bourgain \cite{BB09} later gave a sharper version of Kahane's construction. Erd\'elyi \cite{Erd01} resolved a conjecture of Saffari on the phase of ultraflat unimodular polynomials: their normalized angular speed is asymptotically uniformly distributed. For coefficients $\pm1$, Balister, Bollob\'as, Morris, Sahasrabudhe and Tiba \cite{BBMST} proved that flat Littlewood polynomials exist, answering Littlewood's question. Very recently, the manuscript \cite{AMM} showed that for every $\eta>0$ and every sufficiently large $N$ there are signs with $\max_{|z|=1}|\sum_{k<N}\varepsilon_kz^k|\le(1+\eta)\sqrt N$, which answers the analogue of Erd\H{o}s's question for coefficients $\pm1$ negatively; as a consequence, the largest merit factor \cite{Gol82} of binary sequences of length $N$ tends to infinity.

Flatness up to constant factors does not give $L^1$ norms close to the $L^2$ norm. For instance, the normalized Rudin--Shapiro polynomials satisfy $\|P\|_\infty\le\sqrt2\,\|P\|_2$, but their values are asymptotically uniformly distributed in the disc of radius $\sqrt2$ \cite{Rod17}, so that $\|P\|_1/\|P\|_2\to2\sqrt2/3$. What is needed here is an upper bound of the form $\|R\|_\infty\le(1+o(1))\|R\|_2$ for the balanced part $R$ of a $0$--$1$ polynomial.

\subsection{Results}

\begin{theorem}\label{thm:negative}
$\lim_{n\to\infty}\beta_n=1$. In particular, there is no $\varepsilon>0$ for which \eqref{eq:bourgain} holds.
\end{theorem}

Theorem~\ref{thm:negative} follows from the following biased version of the main theorem of \cite{AMM}.

\begin{theorem}\label{thm:biased}
For every $\rho\in(0,1)$ and $\eta>0$ there is $N_0$ such that for every integer $N\ge N_0$ there is a set $S\subset\{0,1,\dots,N-1\}$ with
\begin{equation}\label{eq:biased}
\max_{t\in\T}\Big|\sum_{k=0}^{N-1}\big(1_S(k)-\rho\big)\e(kt)\Big|\le(1+\eta)\sqrt{\rho(1-\rho)N}.
\end{equation}
In particular $\big||S|-\rho N\big|\le(1+\eta)\sqrt{\rho(1-\rho)N}$.
\end{theorem}

The constant $1$ in \eqref{eq:biased} cannot be improved: by Parseval, $R=\sum_{k<N}(1_S(k)-\rho)\e(k\,\cdot)$ satisfies $\|R\|_2^2=\rho(1-\rho)N+(|S|-\rho N)(1-2\rho)$, and $|S|-\rho N=R(0)=O(\sqrt N)$ under \eqref{eq:biased}. For $\rho=\frac12$, writing $\varepsilon_k=2\cdot1_S(k)-1$, Theorem~\ref{thm:biased} is the main theorem of \cite{AMM}. The bound \eqref{eq:biased} is one-sided; together with Parseval it gives $\|R\|_1\ge(1-o(1))\|R\|_2$ as $\eta\to0$, and this is all that is needed. Since $N_0(\rho,\eta)$ may be taken nonincreasing in $\eta$, choosing $\eta=\eta_N\to0$ slowly enough that $N\ge N_0(\rho,\eta_N)$ gives sets with $\max_t|\sum_{k<N}(1_S(k)-\rho)\e(kt)|\le(1+o(1))\sqrt{\rho(1-\rho)N}$ for every large $N$.

For sets of Theorem~\ref{thm:biased} we obtain $\|X_S\|_1\ge\sqrt{1-\rho}/(1+\eta)-o(1)$ (Section~\ref{sec:deduction}), hence $\|X_S\|_1\to\sqrt{1-\rho}$ along the sets with $\eta=\eta_N$. This is optimal: every $S\subset\{0,\dots,N-1\}$ of density $\rho_S=|S|/N$ satisfies $\|X_S\|_1\le\sqrt{1-\rho_S}+o(1)$ (Proposition~\ref{prop:densitycap}). In particular, sets with $\|X_S\|_1\to1$ must have density tending to $0$ in every interval containing them.

All the main ideas of the construction are taken from \cite{AMM}: an iteratively defined model polynomial, the spreading of its Fourier coefficients by quadratic phases, a packing of signed intervals, sampling along quadratic paths with at most one stationary-phase contribution at each angle, and a rounding procedure with a small defect. We use four lemmas of \cite{AMM} as stated (Section~\ref{sec:tools}). The remaining steps of \cite{AMM} are written for the value interval $[-1,1]$, and we redo them for $[-\rho,1-\rho]$. Four modifications are needed:
\begin{enumerate}
\item[(M1)] the model polynomial is the iteration of \cite{AMM}, written for $u=(1+p)/2$, started at $u_0=\rho$ instead of $u_0=\frac12$ (Lemma~\ref{lem:model});
\item[(M2)] the truncation of the spread polynomial is carried out with tolerance $\delta\min(\rho,1-\rho)$, so that the auxiliary polynomial takes values in $[-\rho,1-\rho]$ (Proposition~\ref{prop:aux});
\item[(M3)] the velocity of the quadratic paths is divided by the energy $M$ of the model polynomial, so that the maximum of the sampled polynomial is $\approx\sqrt{MN}$ rather than $\approx\sqrt N$ (Proposition~\ref{prop:aux});
\item[(M4)] the passage from energy to rounding defect uses the two-point defect $\sum_kp_k(1-p_k)$ and the smallness of the mean of the sampled vector (Section~\ref{sec:rounding}).
\end{enumerate}
Apart from Lemmas~\ref{lem:rounding} and~\ref{lem:packing}, which we quote from \cite{AMM}, we write out all steps in full; for completeness, the two stationary-phase lemmas of \cite{AMM} that we use (Lemmas~\ref{lem:integral} and~\ref{lem:sums}) are proved in Appendix~\ref{app:stationary}. The threshold $N_0(\rho,\eta)$ is not effective, so Theorem~\ref{thm:negative} gives no rate for $1-\beta_n$; Bourgain's bound $1-\beta_n\ge c\log n/n$ is the only general lower bound we know.

The paper is organised as follows. Section~\ref{sec:deduction} deduces Theorem~\ref{thm:negative} from Theorem~\ref{thm:biased}. Section~\ref{sec:tools} states the four lemmas of \cite{AMM}. Sections~\ref{sec:model}--\ref{sec:rounding} prove Theorem~\ref{thm:biased}, and Appendix~\ref{app:stationary} proves Lemmas~\ref{lem:integral} and~\ref{lem:sums}.

\section{From biased flatness to \texorpdfstring{$L^1$}{L1} norms}\label{sec:deduction}
Write $\cD_N(t)=\sum_{k=0}^{N-1}\e(kt)$ and $\|t\|$ for the distance from $t$ to $\Z$.

\begin{proposition}[density cap]\label{prop:densitycap}
Let $S\subset\{0,\dots,N-1\}$, $n=|S|\ge1$ and $\rho_S=n/N$. Then $\|X_S\|_1\le\sqrt{1-\rho_S}+\sqrt n\,(1+\log N)/N$.
\end{proposition}

\begin{proof}
Write $P_S=R_S+\rho_S\cD_N$ with $R_S=\sum_{k<N}(1_S(k)-\rho_S)\e(k\,\cdot)$. By Parseval, $\|R_S\|_2^2=n(1-\rho_S)^2+(N-n)\rho_S^2=\rho_S(1-\rho_S)N$. Since $|\cD_N(t)|=|\sin\pi Nt/\sin\pi t|\le\min(N,1/(2\|t\|))$,
\[
\|\cD_N\|_1\le2\int_0^{1/2}\min\Big(N,\frac1{2t}\Big)dt=1+\log N .
\]
Hence $\|P_S\|_1\le\|R_S\|_2+\rho_S\|\cD_N\|_1\le\sqrt{\rho_S(1-\rho_S)N}+\rho_S(1+\log N)$, and we divide by $\sqrt n=\sqrt{\rho_SN}$.
\end{proof}

\begin{proof}[Proof of Theorem~\ref{thm:negative}, assuming Theorem~\ref{thm:biased}]
Let $\rho\in(0,1)$ and $\eta>0$, let $N\ge N_0(\rho,\eta)$, and let $S$ be given by Theorem~\ref{thm:biased}. Put $E=(1+\eta)\sqrt{\rho(1-\rho)N}$ and $R=\sum_{k<N}(1_S(k)-\rho)\e(k\,\cdot)$, so that $\|R\|_\infty\le E$ and $\big||S|-\rho N\big|\le E$. By orthogonality,
\[
|S|(1-\rho)=\sum_{k\in S}\big(1_S(k)-\rho\big)=\int_\T P_S\,\overline R\le\|R\|_\infty\|P_S\|_1 .
\]
Hence $\|P_S\|_1\ge|S|(1-\rho)/E$, and
\begin{equation}\label{eq:explicit}
\|X_S\|_1\ge\frac{\sqrt{|S|}\,(1-\rho)}{E}\ge\sqrt{1-\frac E{\rho N}}\;\frac{\sqrt{1-\rho}}{1+\eta}\ \longrightarrow\ \frac{\sqrt{1-\rho}}{1+\eta}\qquad(N\to\infty),
\end{equation}
where we used $|S|\ge\rho N-E\ge0$, valid for large $N$.

Now let $\varepsilon>0$, and fix $\rho$ and $\eta$ with $\sqrt{1-\rho}/(1+\eta)>1-\varepsilon/2$. Given a large integer $n$, let $N=\lfloor n/\rho\rfloor$, so that $0\le n-\rho N<\rho$, and let $S$ be as above. Then $m:=\big|n-|S|\big|\le E+1$, and for large $n$ both $|S|\ge\rho N-E\ge m$ and $N-|S|\ge(1-\rho)N-E\ge m$. Adding $m$ elements of $\{0,\dots,N-1\}\setminus S$ to $S$, or deleting $m$ elements of $S$, we obtain a set $S'$ with $|S'|=n$ and $P_{S'}=P_S\pm P_T$ for a set $T$ with $|T|=m$. Since $\|P_T\|_1\le\|P_T\|_2=\sqrt m$,
\[
\|X_{S'}\|_1\ge\frac{\|P_S\|_1-\sqrt m}{\sqrt n}\ge\frac{(\rho N-E)(1-\rho)/E-\sqrt{E+1}}{\sqrt{\rho N+\rho}}\ \longrightarrow\ \frac{\sqrt{1-\rho}}{1+\eta}>1-\frac\varepsilon2\qquad(n\to\infty).
\]
Hence $\beta_n>1-\varepsilon$ for all large $n$. Since $\beta_n\le1$, we conclude that $\beta_n\to1$.
\end{proof}

\begin{remark}[autocorrelations]
Let $g=|X_S|^2-1=n^{-1}\sum_{h\ne0}r(h)\e(h\theta)$ and $Y=|X_S|$. Since $\int Y^2=1$, we have $1-\|X_S\|_1=\frac12\int(Y-1)^2$. As $(Y-1)^2\le|Y-1|(Y+1)=|g|$, and as $\int|g|=\int|Y-1|(Y+1)\le(\int(Y-1)^2)^{1/2}(\int(Y+1)^2)^{1/2}\le2\,(\int(Y-1)^2)^{1/2}$ by Cauchy--Schwarz,
\begin{equation}\label{eq:autocorr}
\tfrac18\|g\|_1^2\le1-\|X_S\|_1\le\tfrac12\|g\|_1 .
\end{equation}
Hence Theorem~\ref{thm:negative} gives sets $S$ with $\big\|\sum_{h\ne0}r(h)\e(h\,\cdot)\big\|_{L^1(\T)}=o(|S|)$.
\end{remark}

By \eqref{eq:explicit} and Proposition~\ref{prop:densitycap}, the sets of Theorem~\ref{thm:biased} have asymptotically the largest possible normalized $L^1$ norm among subsets of $\{0,\dots,N-1\}$ of density $\rho$. Moreover $|R|^2/|S|\le(1+\eta)^2(1-\rho)+o(1)$ while $\|R\|_2^2/|S|\to1-\rho$, so $|R|^2/|S|$ is close to the constant $1-\rho$ in $L^1(\T)$ when $\eta$ is small: the balanced part is flat, and the only deviation of $|X_S|$ from $1$ comes from the peak $\rho\cD_N/\sqrt{|S|}$ and from the factor $\sqrt{1-\rho}$.

\section{Four lemmas}\label{sec:tools}
We use the following four results of \cite{AMM}; none of them involves the values of the auxiliary polynomial. For $X=(X_0,\dots,X_{N-1})\in\R^N$ write $Q_X(t)=\sum_{k=0}^{N-1}X_k\e(kt)$, and for $X\in[-1,1]^N$ put $\mu(X)=\frac12\sum_{k<N}(1-|X_k|)$. For $\xi\in\T$ and $0<\ell<1$ let $I(\xi,\ell)=\{\xi+x\bmod1:\ |x|\le\ell/2\}$, the closed arc with centre $\xi$ and length $\ell$.

\begin{lemma}[{\cite[Lemma~2.2]{AMM}}]\label{lem:rounding}
There is an absolute constant $C$ such that for every $N\ge1$ and every $X\in[-1,1]^N$ there are signs $\varepsilon_0,\dots,\varepsilon_{N-1}\in\{-1,1\}$ with
\[
\max_{t\in\T}\Big|\sum_{k=0}^{N-1}(\varepsilon_k-X_k)\e(kt)\Big|\le C\Big(1+\sqrt{\mu(X)\log\big(80N/\mu(X)\big)}\Big),
\]
where the square-root term is $0$ if $\mu(X)=0$.
\end{lemma}

The proof of Lemma~\ref{lem:rounding} in \cite{AMM} is a dyadic rounding argument based on the partial colouring method of Beck \cite{Bec81} and Spencer \cite{Spe85}, in the constructive form of Lovett and Meka \cite{LM15}.

\begin{lemma}[{\cite[Lemma~3.1]{AMM}}]\label{lem:packing}
Let $m\ge1$ and $r\ge2$ be integers, let $a_1,\dots,a_r\in\Z^m\setminus\{0\}$ be pairwise nonparallel, and let $w_1,\dots,w_r>0$ satisfy $2\sum_iw_i<1$. Then there are an integer $H\ge1$ and points $\theta_1,\dots,\theta_H\in\T^m$ such that the $2rH$ closed arcs $I(\sigma a_i\cdot\theta_h,w_i/H)$, $1\le i\le r$, $1\le h\le H$, $\sigma\in\{-1,1\}$, are pairwise disjoint.
\end{lemma}

The proof of Lemma~\ref{lem:packing} in \cite{AMM} uses the theorem of Pippenger and Spencer \cite{PS89} on the chromatic index of hypergraphs, in the form of Alon and Yuster \cite{AY05}.

\begin{lemma}[{\cite[Lemma~4.2]{AMM}}]\label{lem:integral}
Let $g\in C_c^\infty(\R)$ and $b\in\R\setminus\{0\}$ be fixed. As $T\to\infty$,
\[
\int_\R g(x)\,\e\big(Tbx^2/2-ux\big)\,dx=\frac{\e\big(\mathrm{sgn}(b)/8-u^2/(2Tb)\big)}{\sqrt{T|b|}}\Big(g\Big(\frac u{Tb}\Big)+O_{g,b}(T^{-1})\Big),
\]
uniformly for $u\in\R$.
\end{lemma}

\begin{lemma}[{\cite[Lemma~5.1]{AMM}}]\label{lem:sums}
Let $\chi\in C_c^\infty(\R)$, $x_*\in\R$, $\lambda\in\R\setminus\{0\}$, and let $J\subset\R$ be compact. For $b\in J$ and $\ell\in\Z$ put $x_{b,\ell}=x_*+(\ell-b)/\lambda$ and $\omega_{b,\ell}=(b-\ell)x_*-(b-\ell)^2/(2\lambda)$. As $N\to\infty$ through the integers, uniformly for $b\in J$,
\[
\frac1{\sqrt N}\sum_{k\in\Z}\chi(k/N)\,\e\Big(bk+\frac{N\lambda}2\Big(\frac kN-x_*\Big)^2\Big)=\frac1{\sqrt{|\lambda|}}\sum_{\ell\in\Z}\chi(x_{b,\ell})\,\e\Big(\frac{\mathrm{sgn}\lambda}8+N\omega_{b,\ell}\Big)+O(N^{-1}).
\]
The sum on the right has finitely many nonzero terms, with indices in a fixed finite set for all $b\in J$, and the implied constant depends only on $\chi,x_*,\lambda,J$.
\end{lemma}

\section{The model polynomial}\label{sec:model}
This section contains modification (M1).

\begin{lemma}\label{lem:model}
Let $\rho\in(0,1)$. Put $U_0=\rho$ and, for $j\ge1$ and $y=(y_1,\dots,y_j)\in\T^j$,
\[
U_j(y)=U_{j-1}(y_1,\dots,y_{j-1})+U_{j-1}(y_1,\dots,y_{j-1})\big(1-U_{j-1}(y_1,\dots,y_{j-1})\big)\cos(2\pi y_j).
\]
Let $F_j=U_j-\rho$ and $V_j=\int_{\T^j}F_j^2$. Then for every $j\ge0$:
\begin{enumerate}
\item[(a)] $U_j$ is a real trigonometric polynomial with $0\le U_j\le1$; hence $-\rho\le F_j\le1-\rho$;
\item[(b)] $\int_{\T^j}F_j=0$;
\item[(c)] $V_j$ is nondecreasing in $j$, $V_j\le\rho(1-\rho)$, and $\rho(1-\rho)-V_j\le\big(\frac1{\rho(1-\rho)}+\frac j2\big)^{-1}$;
\item[(d)] if $j\ge1$, every $s\in\Z^j$ with $\widehat{F_j}(s)\ne0$ is nonzero and its last nonzero coordinate equals $1$ or $-1$;
\item[(e)] if $j\ge2$, the Fourier support of $F_j$ is $\{\pm s_1,\dots,\pm s_D\}$ for some $D\ge2$ and pairwise nonparallel $s_1,\dots,s_D\in\Z^j\setminus\{0\}$.
\end{enumerate}
For $\rho=\frac12$, $p_j=2U_j-1$ satisfies $p_j=p_{j-1}+\frac12(1-p_{j-1}^2)\cos(2\pi y_j)$, $p_0=0$, which is the model iteration of \cite{AMM}.
\end{lemma}

\begin{proof}
(a) If $u\in[0,1]$ and $c\in[-1,1]$, then $u+u(1-u)c\in[u-u(1-u),u+u(1-u)]=[u^2,1-(1-u)^2]\subset[0,1]$. By induction, $U_j$ is a polynomial in the functions $\cos2\pi y_i$, $i\le j$, with values in $[0,1]$.

(b) Integrating in $y_j$ gives $\int_\T U_j\,dy_j=U_{j-1}$, since $\int_\T\cos2\pi y_j\,dy_j=0$. By Fubini and induction, $\int U_j=\rho$.

(c) Integrating in $y_j$ and using $\int_\T\cos^22\pi y\,dy=\frac12$,
\[
\int_\T F_j^2\,dy_j=F_{j-1}^2+\tfrac12U_{j-1}^2(1-U_{j-1})^2 .
\]
By (b), $\int U_{j-1}(1-U_{j-1})=\rho-\int U_{j-1}^2=\rho-\rho^2-V_{j-1}$, so by the Cauchy--Schwarz inequality
\[
V_j=V_{j-1}+\tfrac12\int U_{j-1}^2(1-U_{j-1})^2\ge V_{j-1}+\tfrac12\big(\rho(1-\rho)-V_{j-1}\big)^2 .
\]
On the other hand $V_j=\int U_j^2-\rho^2\le\int U_j-\rho^2=\rho(1-\rho)$, as $U_j^2\le U_j$. Thus $e_j=\rho(1-\rho)-V_j\in[0,\frac14]$ satisfies $e_j\le e_{j-1}(1-e_{j-1}/2)$. If $e_j>0$ this gives $1/e_j\ge(1/e_{j-1})(1+e_{j-1}/2)=1/e_{j-1}+\frac12$, and (c) follows by induction from $e_0=\rho(1-\rho)$.

(d) Let $j\ge1$ and $f=U_{j-1}(1-U_{j-1})$, a real trigonometric polynomial on $\T^{j-1}$. Since $\cos2\pi y_j=\frac12(\e(y_j)+\e(-y_j))$, for $(s',\sigma)\in\Z^{j-1}\times\Z$ we have
\[
\widehat{U_j}(s',0)=\widehat{U_{j-1}}(s'),\qquad\widehat{U_j}(s',\pm1)=\tfrac12\widehat f(s'),\qquad\widehat{U_j}(s',\sigma)=0\ \ (|\sigma|\ge2).
\]
Hence the Fourier support of $F_j$ is contained in $(\mathrm{supp}\,\widehat{F_{j-1}}\times\{0\})\cup(\Z^{j-1}\times\{-1,1\})$, and $\widehat{F_j}(0)=0$ by (b). Since $F_1=\rho(1-\rho)\cos2\pi y_1$, (d) follows by induction on $j$.

(e) Let $j\ge2$. If $s$ and $s'$ lie in the Fourier support and $s=ts'$ with $t\in\R$, then $s$ and $s'$ have their last nonzero coordinates in the same position, and both equal $\pm1$ by (d); so $t=\pm1$. Since $F_j$ is real, its Fourier support is symmetric, and choosing one element from each pair $\{s,-s\}$ gives pairwise nonparallel $s_1,\dots,s_D$. Finally, $\widehat{F_j}(e_1)=\widehat{F_1}(1)=\rho(1-\rho)/2\ne0$, and $\widehat{F_j}(e_2)=\widehat{F_2}(0,1)=\frac12\int U_1(1-U_1)=\frac12(\rho(1-\rho)-V_1)>0$, because $V_1=\rho^2(1-\rho)^2/2<\rho(1-\rho)$. As $e_1$ and $e_2$ are not parallel, $D\ge2$. The last assertion is a direct substitution.
\end{proof}

\section{The auxiliary polynomial}\label{sec:aux}
Let $K_\delta=\big((1+\delta)^3/(1-\delta)\big)^{1/2}$ and $m_\rho=\min(\rho,1-\rho)$. This section contains modifications (M2) and (M3).

\begin{proposition}\label{prop:aux}
Let $\rho\in(0,1)$ and $0<\delta<\frac1{20}$. There are integers $m\ge1$ and $r\ge2$, pairwise nonparallel vectors $a_1,\dots,a_r\in\Z^m\setminus\{0\}$, a real trigonometric polynomial $F$ on $\T^m$, a number $M$ and a vector $v\in\R^m$ such that, with $A=\{\pm a_1,\dots,\pm a_r\}$ and $\lambda_a=a\cdot v$:
\begin{enumerate}
\item[(i)] $(1-\delta)\rho(1-\rho)\le M\le\rho(1-\rho)$;
\item[(ii)] $\widehat F(a)=0$ for $a\notin A$;
\item[(iii)] $-\rho\le F\le1-\rho$ on $\T^m$, and $\|F\|_2\ge(1-3\delta)\sqrt M$;
\item[(iv)] $\lambda_a\ne0$ for $a\in A$, and $\sum_{a\in A}|\lambda_a|\le1/(1+\delta)$;
\item[(v)] $|\widehat F(a)|\le K_\delta\sqrt{M|\lambda_a|}$ for $a\in A$.
\end{enumerate}
\end{proposition}

\begin{proof}
\emph{Step 0: the model.} By Lemma~\ref{lem:model}(c) there is $d\ge2$ with $M:=V_d\ge(1-\delta)\rho(1-\rho)$; also $M\le\rho(1-\rho)$, which is (i). Write $q=F_d=\sum_{s\in\cS}c_s\e(s\cdot y)$, $y\in\T^d$, where $\cS=\{\pm s_1,\dots,\pm s_D\}$ is the Fourier support of $q$ as in Lemma~\ref{lem:model}(e). Then $c_s\ne0$ and $c_{-s}=\overline{c_s}$ for $s\in\cS$, $D\ge2$, the $s_i$ are pairwise nonparallel, $-\rho\le q\le1-\rho$, and $\sum_{s\in\cS}|c_s|^2=M$ by Parseval.

\emph{Step 1: curvatures.} Choose $\gamma\in\R^d$ with $s\cdot\gamma\ne0$ for all $s\in\cS$ (this excludes finitely many hyperplanes), and put $\nu_s=|s\cdot\gamma|>0$; note $\nu_{-s}=\nu_s$. For each $j\in\{1,\dots,D\}$ choose $b_j^0\in s_j^\perp$ with $s_l\cdot b_j^0\ne0$ for all $l\ne j$. This is possible: $s_j^\perp$ has dimension $d-1\ge1$, and for $l\ne j$ the subspace $s_j^\perp\cap s_l^\perp$ is a proper subspace of $s_j^\perp$, because $s_l$ is not parallel to $s_j$; a finite union of proper subspaces does not cover $s_j^\perp$. Choose $b_j^1\in\R^d$ with
\[
s_j\cdot b_j^1=\frac{|c_{s_j}|^2}{\nu_{s_j}\prod_{l\ne j}|s_j\cdot b_l^0|}>0,
\]
and for $\tau>0$ put $b_j=b_j^0+\tau b_j^1$ and $\Delta_s=\prod_{l=1}^D|s\cdot b_l|$. For $s=\pm s_j$ we have $|s\cdot b_j|=\tau\,s_j\cdot b_j^1$ and $|s\cdot b_l|\to|s_j\cdot b_l^0|>0$ as $\tau\to0$ for $l\ne j$, so $\Delta_s/\tau\to|c_s|^2/\nu_s$. Fix $\tau>0$ so small that $s\cdot b_l\ne0$ for all $s\in\cS$ and all $l$, and
\begin{equation}\label{eq:curv}
(1-\delta)\,\tau|c_s|^2/\nu_s\le\Delta_s\le(1+\delta)\,\tau|c_s|^2/\nu_s\qquad(s\in\cS).
\end{equation}

\emph{Step 2: spreading.} Choose $g\in C_c^\infty((-\frac12,\frac12))$ with $0\le g\le1$ and $(\int_\R g^2)^D\ge1-\delta$, and fix $a_g<\frac12$ with $\mathrm{supp}\,g\subset[-a_g,a_g]$. For $B>0$, $y\in\T^d$ and $z\in[-\frac12,\frac12]^D$ put
\[
G(y,z)=\Big(\prod_{j=1}^Dg(z_j)\Big)\,q\Big(y+\frac B2\sum_{j=1}^Db_jz_j^2\Big).
\]
Since $g$ vanishes near $\pm\frac12$, $G$ extends to a smooth function on $\T^{d+D}$. As $\prod_jg(z_j)\in[0,1]$ and $0\in[-\rho,1-\rho]$, we have $-\rho\le G\le1-\rho$; and by translation invariance in $y$, $\|G\|_2^2=M(\int g^2)^D\ge(1-\delta)M$. Integrating first in $y$, for $(s',u)\in\Z^d\times\Z^D$,
\begin{equation}\label{eq:Ghat}
\widehat G(s',u)=0\ \ (s'\notin\cS),\qquad\widehat G(s,u)=c_s\prod_{j=1}^DI_{s,j}(u_j),\qquad I_{s,j}(w)=\int_\R g(x)\,\e\Big(\frac{B(s\cdot b_j)}2x^2-wx\Big)dx .
\end{equation}

\emph{Coefficient bound.} By Lemma~\ref{lem:integral} with $T=B$ and $b=s\cdot b_j\ne0$, and since $0\le g\le1$, there are constants $C_{s,j}$ and $B_{s,j}$ with $|I_{s,j}(w)|\le(1+C_{s,j}/B)/\sqrt{B|s\cdot b_j|}$ for all $w\in\R$ and all $B\ge B_{s,j}$. As there are finitely many pairs $(s,j)$, there is $B_1$ such that for $B\ge B_1$
\begin{equation}\label{eq:coef}
|\widehat G(s,u)|\le(1+\delta)\,|c_s|/\sqrt{B^D\Delta_s}\qquad(s\in\cS,\ u\in\Z^D).
\end{equation}

\emph{Boxes.} Let $U_s=\{u\in\Z^D:\ |u_j|\le B|s\cdot b_j|/2\text{ for }1\le j\le D\}$, so that $U_{-s}=U_s=-U_s$. Then
\[
|U_s|=\prod_{j=1}^D\big(2\lfloor B|s\cdot b_j|/2\rfloor+1\big)\le B^D\Delta_s\prod_{j=1}^D\Big(1+\frac1{B|s\cdot b_j|}\Big)\le(1+\delta)B^D\Delta_s
\]
for $B\ge B_2$.

\emph{Tails.} Fix $b\ne0$ and an integer $l\ge1$, and put $I_b(w)=\int_\R g(x)\,\e(Bbx^2/2-wx)\,dx$, so that $I_{s,j}=I_{s\cdot b_j}$. For $|w|>B|b|/2$ and $x\in\mathrm{supp}\,g$, the phase $\phi(x)=Bbx^2/2-wx$ satisfies $|\phi'(x)|=|Bbx-w|\ge|w|-B|b|a_g\ge(1-2a_g)|w|$, and $B<2|w|/|b|$ gives $|w|\ge(B+|w|)/(1+2/|b|)$. Hence $|\phi'|\ge c_b(B+|w|)$ on $\mathrm{supp}\,g$, with $c_b=(1-2a_g)/(1+2/|b|)>0$. Integrating by parts $l$ times with the operator $f\mapsto-\frac d{dx}\big(f/(2\pi i\phi')\big)$ (there are no boundary terms) and using $\phi''=Bb$, $\phi'''=0$, the integral $I_b(w)$ becomes a finite sum of terms of the form $c\int g^{(l-i)}(Bb)^i(\phi')^{-l-i}\e(\phi)$, $0\le i\le l$. Since $B|b|\le|b|(B+|w|)$, each term is $O_{g,b,l}((B+|w|)^{-l})$. So
\[
|I_b(w)|\le C_{g,b,l}(B+|w|)^{-l}\qquad(|w|>B|b|/2,\ B\ge1).
\]
Summing over integers $w$ with $l=D+1\ge3$ gives $\sum_{|w|>B|b|/2}|I_b(w)|=O(B^{-D})$, while trivially $\sum_{w\in\Z}|I_b(w)|\le(B|b|+1)\|g\|_1+O(B^{-D})=O(B)$. Since $\Z^D\setminus U_s$ is covered by the $D$ sets $\{u:|u_j|>B|s\cdot b_j|/2\}$, \eqref{eq:Ghat} gives
\[
\sum_{u\notin U_s}|\widehat G(s,u)|\le|c_s|\sum_{j=1}^D\Big(\sum_{|w|>B|s\cdot b_j|/2}|I_{s,j}(w)|\Big)\prod_{j'\ne j}\Big(\sum_{w\in\Z}|I_{s,j'}(w)|\Big)=O(B^{-1}).
\]
Hence there is $B_3$ such that for $B\ge B_3$
\begin{equation}\label{eq:tail}
\sum_{s\in\cS}\sum_{u\notin U_s}|\widehat G(s,u)|\le\delta\,m_\rho .
\end{equation}

\emph{Step 3: truncation and normalisation} (modification (M2)). Fix $B\ge\max(1,B_1,B_2,B_3)$ and put
\[
T_B(y,z)=\sum_{s\in\cS}\sum_{u\in U_s}\widehat G(s,u)\,\e(s\cdot y+u\cdot z),\qquad F=\frac{T_B}{1+\delta}.
\]
Since $G$ is real, $\widehat G(-s,-u)=\overline{\widehat G(s,u)}$, and since the index set is symmetric under $(s,u)\mapsto(-s,-u)$, $T_B$ and $F$ are real. The Fourier series of the smooth function $G$ converges absolutely, so \eqref{eq:tail} gives $\|G-T_B\|_\infty\le\delta m_\rho$. Therefore
\[
-\rho(1+\delta)\le-\rho-\delta m_\rho\le T_B\le1-\rho+\delta m_\rho\le(1-\rho)(1+\delta),
\]
that is, $-\rho\le F\le1-\rho$. For the energy, $m_\rho^2\le\rho(1-\rho)\le M/(1-\delta)$, so
\[
\|T_B\|_2\ge\|G\|_2-\|G-T_B\|_\infty\ge\sqrt{(1-\delta)M}-\delta\sqrt{\frac M{1-\delta}}=\frac{1-2\delta}{\sqrt{1-\delta}}\sqrt M\ge(1-2\delta)\sqrt M,
\]
and $\|F\|_2\ge\frac{1-2\delta}{1+\delta}\sqrt M\ge(1-3\delta)\sqrt M$. This proves (iii).

\emph{Step 4: support and velocity} (modification (M3)). Let $m=d+D$ and $A=\{(s,u):\ s\in\cS,\ u\in U_s\}\subset\Z^m$. Then $\widehat F$ vanishes off $A$, $A=-A$ and $0\notin A$. If $(s,u)=t(s',u')$ with $(s,u),(s',u')\in A$ and $t\in\R$, then $s=ts'$, so $t=\pm1$ by Lemma~\ref{lem:model}(e). Hence $A=\{\pm a_1,\dots,\pm a_r\}$ with $a_1,\dots,a_r$ pairwise nonparallel, and $r\ge D\ge2$ because $(s_i,0)\in A$ for every $i$. This proves (ii). Put
\[
v=\Big(\frac{\gamma}{\tau B^D(1+\delta)^3M},\,0\Big)\in\R^d\times\R^D .
\]
(In \cite{AMM}, where the model has energy close to $1$, the factor $1/M$ is absent.) For $a=(s,u)\in A$ we get $\lambda_a=s\cdot\gamma/(\tau B^D(1+\delta)^3M)$, so $|\lambda_a|=\nu_s/(\tau B^D(1+\delta)^3M)>0$. By the bound on $|U_s|$, by \eqref{eq:curv} and by $\sum_s|c_s|^2=M$,
\[
\sum_{a\in A}|\lambda_a|=\sum_{s\in\cS}\frac{|U_s|\,\nu_s}{\tau B^D(1+\delta)^3M}\le\sum_{s\in\cS}\frac{\Delta_s\nu_s}{\tau(1+\delta)^2M}\le\sum_{s\in\cS}\frac{|c_s|^2}{(1+\delta)M}=\frac1{1+\delta},
\]
which is (iv). Finally, for $a=(s,u)\in A$, \eqref{eq:coef} gives $|\widehat F(a)|^2=|\widehat G(s,u)|^2/(1+\delta)^2\le|c_s|^2/(B^D\Delta_s)$, and by the lower bound in \eqref{eq:curv}
\[
\frac{|\widehat F(a)|^2}{|\lambda_a|}\le\frac{|c_s|^2\,\tau(1+\delta)^3M}{\Delta_s\nu_s}\le\frac{(1+\delta)^3}{1-\delta}M=K_\delta^2M,
\]
which is (v).
\end{proof}

\begin{remark}
Without the factor $1/M$ in $v$ one still has $\sum|\lambda_a|<1$, but only $|\widehat F(a)|\le K_\delta\sqrt{|\lambda_a|}$; this leads to the bound $K_\delta\sqrt N$ in Proposition~\ref{prop:sampling} instead of $K_\delta\sqrt{MN}$, which is not sufficient.
\end{remark}

\section{Sampling}\label{sec:sampling}

\begin{proposition}\label{prop:sampling}
Let $\rho\in(0,1)$, $0<\delta<\frac1{20}$, and let $F,A,M,v$ be as in Proposition~\ref{prop:aux}. There are vectors $X_N=(X_{N,0},\dots,X_{N,N-1})\in[-\rho,1-\rho]^N$, $N\ge1$, such that
\[
\limsup_{N\to\infty}\frac{\|Q_{X_N}\|_\infty}{\sqrt N}\le K_\delta\sqrt M\qquad\text{and}\qquad\lim_{N\to\infty}\frac1N\sum_{k=0}^{N-1}X_{N,k}^2\ge(1-7\delta)M .
\]
\end{proposition}

\begin{proof}
\emph{Packing.} Apply Lemma~\ref{lem:packing} to $a_1,\dots,a_r$ with $w_i=|\lambda_{a_i}|>0$. Since $|\lambda_{-a}|=|\lambda_a|$, we have $2\sum_iw_i=\sum_{a\in A}|\lambda_a|\le1/(1+\delta)<1$. We obtain $H\ge1$ and $\theta_1,\dots,\theta_H\in\T^m$ such that the closed arcs
\[
I_{a,h}=I\big(a\cdot\theta_h,|\lambda_a|/H\big),\qquad a\in A,\ 1\le h\le H,
\]
are pairwise disjoint. Each has length $|\lambda_a|/H<1$. Fix real lifts $\tilde\theta_h\in\R^m$ of the $\theta_h$.

\emph{Cutoffs and definition.} Let $x_h=(h-\frac12)/H$, and choose $\chi_h\in C_c^\infty\big(((h-1)/H,h/H)\big)$ with $0\le\chi_h\le1$ and $\sum_{h=1}^H\int_\R\chi_h^2\ge1-\delta$. All data are now fixed, independently of $N$. For $N\ge1$ and $0\le k<N$ put
\[
X_{N,k}=\sum_{h=1}^H\chi_h(k/N)\,F\Big(k\theta_h+\frac N2\Big(\frac kN-x_h\Big)^2v\Big).
\]
This is real and does not depend on the choice of lifts.

\emph{Range.} The supports of the $\chi_h$ are pairwise disjoint, so at most one term is nonzero, and $X_{N,k}=c\,F(\cdot)$ with $c\in[0,1]$. Since $[-\rho,1-\rho]$ is convex and contains $0$, Proposition~\ref{prop:aux}(iii) gives $X_{N,k}\in[-\rho,1-\rho]$. As each $\chi_h$ is supported in $(0,1)$, we may extend the sum over $k$ in $Q_{X_N}$ to all $k\in\Z$.

\emph{Maximum.} Writing $F(y)=\sum_{a\in A}\widehat F(a)\e(a\cdot y)$, for $t\in[0,1]$
\[
\frac{Q_{X_N}(t)}{\sqrt N}=\sum_{h=1}^H\sum_{a\in A}\widehat F(a)\,\frac1{\sqrt N}\sum_{k\in\Z}\chi_h(k/N)\,\e\Big(k(t+a\cdot\tilde\theta_h)+\frac{N\lambda_a}2\Big(\frac kN-x_h\Big)^2\Big).
\]
Apply Lemma~\ref{lem:sums} to each of the finitely many pairs $(a,h)$, with $\chi=\chi_h$, $x_*=x_h$, $\lambda=\lambda_a$ and $b=t+a\cdot\tilde\theta_h\in J_{a,h}=[a\cdot\tilde\theta_h,a\cdot\tilde\theta_h+1]$. The term with index $\ell$ is nonzero only if $x_h+(\ell-b)/\lambda_a\in\mathrm{supp}\,\chi_h\subset(x_h-\frac1{2H},x_h+\frac1{2H})$, that is, only if $|\ell-b|<|\lambda_a|/(2H)$. This means that $-t\bmod1$ lies in the interior of $I_{a,h}$. By disjointness, for each $t$ at most one pair $(a,h)$ has a nonzero main term, and for that pair at most one $\ell$ does, since $|\lambda_a|/H<1$. That term has modulus at most $|\widehat F(a)|/\sqrt{|\lambda_a|}\le K_\delta\sqrt M$, by Proposition~\ref{prop:aux}(v) and $0\le\chi_h\le1$. The error terms are $O(N^{-1})$ uniformly in $t$. Hence $\|Q_{X_N}\|_\infty\le(K_\delta\sqrt M+O(N^{-1}))\sqrt N$.

\emph{Energy.} Products of terms with different $h$ vanish, so
\[
\frac1N\sum_kX_{N,k}^2=\sum_{h=1}^H\sum_{a,a'\in A}\widehat F(a)\overline{\widehat F(a')}\;\frac1N\sum_{k\in\Z}\chi_h(k/N)^2\,\e\Big(k\vartheta+\frac N2(\lambda_a-\lambda_{a'})\Big(\frac kN-x_h\Big)^2\Big),
\]
where $\vartheta=(a-a')\cdot\tilde\theta_h$. If $a=a'$, the inner average is a Riemann sum tending to $\int\chi_h^2$. If $\lambda_a\ne\lambda_{a'}$, Lemma~\ref{lem:sums} with amplitude $\chi_h^2$, curvature $\lambda_a-\lambda_{a'}$ and $J=\{\vartheta\}$ shows that the inner sum is $O(\sqrt N)$, so the average is $O(N^{-1/2})$. If $a\ne a'$ but $\lambda_a=\lambda_{a'}$, then $a\cdot\theta_h\ne a'\cdot\theta_h$ in $\T$, because $I_{a,h}\cap I_{a',h}=\emptyset$; so $\vartheta\notin\Z$, and summation by parts with $|\sum_{k=k_1}^{k_2}\e(k\vartheta)|\le2/|1-\e(\vartheta)|$ shows that the inner sum is $O(1)$. By Parseval on $\T^m$,
\[
\lim_{N\to\infty}\frac1N\sum_kX_{N,k}^2=\Big(\sum_h\int\chi_h^2\Big)\|F\|_2^2\ge(1-\delta)(1-3\delta)^2M\ge(1-7\delta)M .\qedhere
\]
\end{proof}

\section{Rounding: proof of Theorem~\ref{thm:biased}}\label{sec:rounding}
This section contains modification (M4). Fix $\rho\in(0,1)$ and $\eta>0$, and let $C$ be the constant of Lemma~\ref{lem:rounding}. Since $K_\delta\to1$ and $\delta\log(1/\delta)\to0$ as $\delta\to0$, we can fix $\delta\in(0,\frac1{20})$ with
\[
\Phi(\delta,\rho):=K_\delta+\frac C2\sqrt{17\delta\log\frac{80}{17\delta\rho(1-\rho)}}<1+\eta .
\]
Let $X_N$ be the vectors of Proposition~\ref{prop:sampling}, write $X_k=X_{N,k}$, and put
\[
p_k=\rho+X_k\in[0,1],\qquad Y_k=2p_k-1\in[-1,1],\qquad \Delta_N=\frac1N\sum_{k=0}^{N-1}p_k(1-p_k)\ge0 .
\]
Since $p(1-p)=\rho(1-\rho)+(1-2\rho)X-X^2$ for $p=\rho+X$,
\[
\Delta_N=\rho(1-\rho)+(1-2\rho)\,\frac1N\sum_kX_k-\frac1N\sum_kX_k^2 .
\]
The mean is small: $|\sum_kX_k|=|Q_{X_N}(0)|=O(\sqrt N)$ by Proposition~\ref{prop:sampling}. Using Proposition~\ref{prop:sampling} and Proposition~\ref{prop:aux}(i),
\[
\limsup_{N\to\infty}\Delta_N\le\rho(1-\rho)-(1-7\delta)(1-\delta)\rho(1-\rho)\le8\delta\rho(1-\rho).
\]
Since $1-|2p-1|=2\min(p,1-p)\le4p(1-p)$ for $p\in[0,1]$,
\[
\mu(Y)=\sum_k\min(p_k,1-p_k)\le2\sum_kp_k(1-p_k)=2N\Delta_N .
\]
Hence $\limsup_N\mu(Y)/N\le16\delta\rho(1-\rho)<u_0$, where $u_0=17\delta\rho(1-\rho)<1$. Lemma~\ref{lem:rounding} applied to $Y$ gives signs $\varepsilon_k$ with $\|Q_{\varepsilon-Y}\|_\infty\le C\big(1+\sqrt{\mu(Y)\log(80N/\mu(Y))}\big)$. Let $S=\{k:\varepsilon_k=1\}$. Since $1_S(k)=\frac{1+\varepsilon_k}2$ and $p_k=\frac{1+Y_k}2$,
\[
1_S(k)-\rho=\frac{\varepsilon_k-Y_k}2+X_k,\qquad\text{hence}\qquad\max_t\Big|\sum_{k<N}(1_S(k)-\rho)\e(kt)\Big|\le\tfrac12\|Q_{\varepsilon-Y}\|_\infty+\|Q_{X_N}\|_\infty .
\]
For large $N$ we have $\mu(Y)\le u_0N$. Since $u\mapsto u\log(80/u)$ is increasing on $(0,1]$, it follows that
\[
\mu(Y)\log\frac{80N}{\mu(Y)}\le Nu_0\log\frac{80}{u_0}
\]
(the left side is read as $0$ if $\mu(Y)=0$). Dividing by $\sqrt{\rho(1-\rho)N}$ and using $M\le\rho(1-\rho)$,
\[
\limsup_{N\to\infty}\frac{\max_t\big|\sum_{k<N}(1_S(k)-\rho)\e(kt)\big|}{\sqrt{\rho(1-\rho)N}}\le\frac C2\sqrt{\frac{u_0\log(80/u_0)}{\rho(1-\rho)}}+K_\delta=\Phi(\delta,\rho)<1+\eta .
\]
This proves \eqref{eq:biased} for all $N\ge N_0$, and the bound on $\big||S|-\rho N\big|$ is the case $t=0$. \qed

\section{Remarks}
(1) The threshold $N_0(\rho,\eta)$ depends on the packing of Lemma~\ref{lem:packing} and on the implied constants in Lemmas~\ref{lem:integral} and~\ref{lem:sums}, and is not effective as stated. Consequently Theorem~\ref{thm:negative} gives no rate for $1-\beta_n$. The true order of $1-\beta_n$ lies between Bourgain's $c\log n/n$ and $o(1)$, and determining it remains open.

(2) The bound \eqref{eq:biased} is one-sided: it gives no lower bound for $|\sum_{k<N}(1_S(k)-\rho)\e(kt)|$, and none is needed.

(3) Since $\|X_{LS}\|_1=\|X_S\|_1$ for every integer $L\ge1$, dilates of the sets of Theorem~\ref{thm:biased} can occur as the polynomials in a rank one construction: dilating by the current height of the tower turns the gaps of $S$ into nonnegative spacers. Whether they can be used to produce rank one transformations with absolutely continuous spectral components requires quantitative information that Theorem~\ref{thm:negative} does not provide; we do not address this question.

\appendix
\section{Proofs of Lemmas~\ref{lem:integral} and~\ref{lem:sums}}\label{app:stationary}
We follow \cite{AMM}. Write $\widehat g(\xi)=\int_\R g(x)\e(-\xi x)\,dx$; for $g\in C_c^\infty(\R)$ this is a Schwartz function.

\begin{proof}[Proof of Lemma~\ref{lem:integral}]
Put $x_0=u/(Tb)$. Completing the square, $Tbx^2/2-ux=Tb(x-x_0)^2/2-u^2/(2Tb)$, so the integral equals $\e(-u^2/(2Tb))\int_\R g(x_0+y)\e(Tby^2/2)\,dy$. For $\epsilon>0$ let $a=\epsilon-iTb$, so that $\Re a>0$. By dominated convergence,
\[
\int_\R g(x_0+y)\,\e(Tby^2/2)\,dy=\lim_{\epsilon\downarrow0}\int_\R g(x_0+y)\,e^{-\pi a y^2}\,dy .
\]
Insert $g(x_0+y)=\int_\R\widehat g(\xi)\e(\xi(x_0+y))\,d\xi$; the resulting double integral converges absolutely, and the Gaussian integral $\int_\R e^{-\pi ay^2+2\pi i\xi y}\,dy=a^{-1/2}e^{-\pi\xi^2/a}$ (principal branch) gives
\[
\int_\R g(x_0+y)\,e^{-\pi ay^2}\,dy=a^{-1/2}\int_\R\widehat g(\xi)\,\e(\xi x_0)\,e^{-\pi\xi^2/a}\,d\xi .
\]
Here $|a^{-1/2}|\le(T|b|)^{-1/2}$ and $|e^{-\pi\xi^2/a}|=e^{-\pi\xi^2\epsilon/|a|^2}\le1$. As $\epsilon\downarrow0$, $a^{-1/2}\to(-iTb)^{-1/2}=\e(\mathrm{sgn}(b)/8)/\sqrt{T|b|}$ and $e^{-\pi\xi^2/a}\to\e(-\xi^2/(2Tb))$, so by dominated convergence
\[
\int_\R g(x)\,\e(Tbx^2/2-ux)\,dx=\frac{\e(\mathrm{sgn}(b)/8-u^2/(2Tb))}{\sqrt{T|b|}}\int_\R\widehat g(\xi)\,\e(\xi x_0)\,\e\Big(-\frac{\xi^2}{2Tb}\Big)d\xi .
\]
Since $\int\widehat g(\xi)\e(\xi x_0)\,d\xi=g(x_0)$ and $|\e(t)-1|\le2\pi|t|$, the last integral differs from $g(x_0)$ by at most $\frac\pi{T|b|}\int_\R\xi^2|\widehat g(\xi)|\,d\xi$, which does not depend on $u$.
\end{proof}

\begin{proof}[Proof of Lemma~\ref{lem:sums}]
For $b\in J$ and $\ell\in\Z$ let $\phi_{b,\ell}(x)=\frac\lambda2(x-x_*)^2+(b-\ell)x$. The function $x\mapsto\chi(x/N)\e(bx+\frac{N\lambda}2(\frac xN-x_*)^2)$ is smooth with compact support, so Poisson summation and the substitution $x=Ny$ give
\[
\frac1{\sqrt N}\sum_{k\in\Z}\chi(k/N)\,\e\Big(bk+\frac{N\lambda}2\Big(\frac kN-x_*\Big)^2\Big)=\sqrt N\sum_{\ell\in\Z}\int_\R\chi(y)\,\e\big(N\phi_{b,\ell}(y)\big)\,dy .
\]
Let $\mathrm{supp}\,\chi\subset[-R,R]$ and $C_0=\sup\{|\lambda(x-x_*)+b|:\ |x|\le R,\ b\in J\}$. For $|\ell|>L:=\max(2C_0,2)$ we have $|\phi'_{b,\ell}|=|\lambda(x-x_*)+b-\ell|\ge|\ell|/2$ on $\mathrm{supp}\,\chi$. Two integrations by parts, using $\phi_{b,\ell}''=\lambda$ and $\phi_{b,\ell}'''=0$, give
\[
\int_\R\chi\,\e(N\phi_{b,\ell})=\frac1{(2\pi iN)^2}\int_\R\Big(\frac{\chi''}{(\phi'_{b,\ell})^2}-\frac{3\lambda\chi'}{(\phi'_{b,\ell})^3}+\frac{3\lambda^2\chi}{(\phi'_{b,\ell})^4}\Big)\e(N\phi_{b,\ell}),
\]
which is $O(N^{-2}\ell^{-2})$ for $|\ell|>L$, uniformly in $b\in J$. These terms contribute $O(N^{-3/2})$ in total. For $|\ell|\le L$, write $N\phi_{b,\ell}(x)=N\lambda x^2/2-ux+N\lambda x_*^2/2$ with $u=N(\ell-b+\lambda x_*)$, and apply Lemma~\ref{lem:integral} with $T=N$, curvature $\lambda$ and this $u$. Since $u/(N\lambda)=x_{b,\ell}$ and $\lambda x_*^2/2-(\ell-b+\lambda x_*)^2/(2\lambda)=\omega_{b,\ell}$, we get
\[
\sqrt N\int_\R\chi\,\e(N\phi_{b,\ell})=\frac1{\sqrt{|\lambda|}}\,\e\Big(\frac{\mathrm{sgn}\lambda}8+N\omega_{b,\ell}\Big)\big(\chi(x_{b,\ell})+O(N^{-1})\big),
\]
uniformly in $b\in J$, by the uniformity in $u$ in Lemma~\ref{lem:integral}. Enlarging $L$ if necessary, $\chi(x_{b,\ell})=0$ for all $|\ell|>L$ and $b\in J$. Summing over $|\ell|\le L$ proves the lemma.
\end{proof}

\begin{thebibliography}{99}

\bibitem{Ais13} C.~Aistleitner, \emph{On a problem of Bourgain concerning the $L^1$-norm of exponential sums}, Math.\ Z.\ \textbf{275} (2013), 681--688.

\bibitem{AH12} C.~Aistleitner and M.~Hofer, \emph{On the maximal spectral type of a class of rank one transformations}, Dyn.\ Syst.\ \textbf{27} (2012), 515--523.

\bibitem{AY05} N.~Alon and R.~Yuster, \emph{On a hypergraph matching problem}, Graphs Combin.\ \textbf{21} (2005), 377--384.

\bibitem{AMM} \emph{Asymptotically minimal maxima of real Littlewood polynomials}, manuscript, September~23, 2026.

\bibitem{BBMST} P.~Balister, B.~Bollob\'as, R.~Morris, J.~Sahasrabudhe and M.~Tiba, \emph{Flat Littlewood polynomials exist}, Ann.\ of Math.\ (2) \textbf{192} (2020), 977--1004.

\bibitem{Bec81} J.~Beck, \emph{Roth's estimate of the discrepancy of integer sequences is nearly sharp}, Combinatorica \textbf{1} (1981), 319--325.

\bibitem{BB09} E.~Bombieri and J.~Bourgain, \emph{On Kahane's ultraflat polynomials}, J.\ Eur.\ Math.\ Soc.\ \textbf{11} (2009), 627--703.

\bibitem{BL01} P.~Borwein and R.~Lockhart, \emph{The expected $L_p$ norm of random polynomials}, Proc.\ Amer.\ Math.\ Soc.\ \textbf{129} (2001), 1463--1472.

\bibitem{Bou93} J.~Bourgain, \emph{On the spectral type of Ornstein's class one transformations}, Israel J.\ Math.\ \textbf{84} (1993), 53--63.

\bibitem{CE15} S.~Choi and T.~Erd\'elyi, \emph{On a problem of Bourgain concerning the $L_p$ norms of exponential sums}, Math.\ Z.\ \textbf{279} (2015), 577--584.

\bibitem{Erd01} T.~Erd\'elyi, \emph{The phase problem of ultraflat unimodular polynomials: the resolution of the conjecture of Saffari}, Math.\ Ann.\ \textbf{321} (2001), 905--924.

\bibitem{Erd20} T.~Erd\'elyi, \emph{Recent progress in the study of polynomials with constrained coefficients}, in: A.~Raigorodskii and M.~Th.~Rassias (eds.), Trigonometric Sums and Their Applications, Springer, Cham, 2020, 29--69.

\bibitem{Erd57} P.~Erd\H{o}s, \emph{Some unsolved problems}, Michigan Math.\ J.\ \textbf{4} (1957), 291--300.

\bibitem{Gol82} M.~J.~E.~Golay, \emph{The merit factor of long low autocorrelation binary sequences}, IEEE Trans.\ Inform.\ Theory \textbf{28} (1982), 543--549.

\bibitem{HL48} G.~H.~Hardy and J.~E.~Littlewood, \emph{A new proof of a theorem on rearrangements}, J.\ London Math.\ Soc.\ \textbf{23} (1948), 163--168.

\bibitem{IS24} P.~Ivanisvili and Y.~Stone, \emph{Sharpening the gap between $L^1$ and $L^2$ norms}, preprint, arXiv:2407.04835 (2024).

\bibitem{Kah80} J.-P.~Kahane, \emph{Sur les polyn\^omes \`a coefficients unimodulaires}, Bull.\ London Math.\ Soc.\ \textbf{12} (1980), 321--342.

\bibitem{Kar98} A.~A.~Karatsuba, \emph{An estimate for the $L_1$-norm of an exponential sum}, Mat.\ Zametki \textbf{64} (1998), 465--468; English transl.\ in Math.\ Notes \textbf{64} (1998), 401--404.

\bibitem{Kle96} I.~Klemes, \emph{The spectral type of the staircase transformation}, Tohoku Math.\ J.\ (2) \textbf{48} (1996), 247--258.

\bibitem{KR97} I.~Klemes and K.~Reinhold, \emph{Rank one transformations with singular spectral type}, Israel J.\ Math.\ \textbf{98} (1997), 1--14.

\bibitem{Kon81} S.~V.~Konyagin, \emph{On the Littlewood problem}, Izv.\ Akad.\ Nauk SSSR Ser.\ Mat.\ \textbf{45} (1981), 243--265; English transl.\ in Math.\ USSR-Izv.\ \textbf{18} (1982), 205--225.

\bibitem{Kor80} T.~W.~K\"orner, \emph{On a polynomial of Byrnes}, Bull.\ London Math.\ Soc.\ \textbf{12} (1980), 219--224.

\bibitem{Lit66} J.~E.~Littlewood, \emph{On polynomials $\sum^n\pm z^m$, $\sum^ne^{\alpha_mi}z^m$, $z=e^{\theta i}$}, J.\ London Math.\ Soc.\ \textbf{41} (1966), 367--376.

\bibitem{Lit68} J.~E.~Littlewood, \emph{Some Problems in Real and Complex Analysis}, D.~C.~Heath, Lexington, MA, 1968.

\bibitem{LM15} S.~Lovett and R.~Meka, \emph{Constructive discrepancy minimization by walking on the edges}, SIAM J.\ Comput.\ \textbf{44} (2015), 1573--1582.

\bibitem{MPS81} O.~C.~McGehee, L.~Pigno and B.~Smith, \emph{Hardy's inequality and the $L^1$ norm of exponential sums}, Ann.\ of Math.\ (2) \textbf{113} (1981), 613--618.

\bibitem{Orn72} D.~S.~Ornstein, \emph{On the root problem in ergodic theory}, in: Proc.\ Sixth Berkeley Sympos.\ Math.\ Statist.\ Probab., Vol.~II, Univ.\ California Press, Berkeley, 1972, 347--356.

\bibitem{PS89} N.~Pippenger and J.~Spencer, \emph{Asymptotic behavior of the chromatic index for hypergraphs}, J.\ Combin.\ Theory Ser.\ A \textbf{51} (1989), 24--42.

\bibitem{Rod17} B.~Rodgers, \emph{On the distribution of Rudin--Shapiro polynomials and lacunary walks on $SU(2)$}, Adv.\ Math.\ \textbf{320} (2017), 993--1008.

\bibitem{Rud59} W.~Rudin, \emph{Some theorems on Fourier coefficients}, Proc.\ Amer.\ Math.\ Soc.\ \textbf{10} (1959), 855--859.

\bibitem{San26} T.~Sanders, \emph{The Wiener norm of lifts of dense Sidon sets}, preprint, arXiv:2609.31335 (2026).

\bibitem{Sha51} H.~S.~Shapiro, \emph{Extremal problems for polynomials and power series}, Master's thesis, Massachusetts Institute of Technology, 1951.

\bibitem{Spe85} J.~Spencer, \emph{Six standard deviations suffice}, Trans.\ Amer.\ Math.\ Soc.\ \textbf{289} (1985), 679--706.

\end{thebibliography}

\bigskip
\noindent\emph{AI disclosure.} This result was produced in interaction with Grok Bot running on Cursor agents.

\end{document}
